\section{Methods}
\subsection{Timing-resolution definition}
CTR is measured from the distribution of timing residuals using method F1 of Rainio
et al.~\cite{rainio2025}. A uniformly binned histogram is represented by its bin centres
and counts. The maximum bin defines the peak height; when several bins share the maximum,
the middlemost maximum is selected. On either side of this bin, the nearest count strictly
below half the peak height brackets a crossing with its adjacent inner bin. Linear
interpolation between the two bin centres locates each crossing, and their separation is
the FWHM. No Gaussian fit or conversion from a standard deviation is used for this measure.

The physical bin width and histogram phase are fixed for a given waveform mode; zero is a
bin centre. The benchmark settings use $10\,\mathrm{ps}$ bins for energy-waveform timing and
$5\,\mathrm{ps}$ bins for timing-waveform timing. LED and ML residuals are evaluated with
the same estimator. The root mean squared error (RMSE), measured relative to the configured
true TOF, is retained as a complementary measure that also reflects residual bias.

\subsection{Control-fitted event selection}
The two-detector photopeak selection is fitted from the control energy amplitudes using
iterative Gaussian peak fits. For timing-mode analysis, accepted threshold-exceeding
intervals must also satisfy control-fitted time-over-threshold (ToT) ranges.
Baseline quality is evaluated in a trigger-relative pre-pulse interval, from
$-2$ to $-1\,\mathrm{ns}$ in the supplied configuration. For each detector and waveform family,
the baseline root mean square noise is compared with a control-derived robust upper limit,
obtained from the median plus five scaled median absolute deviations. Events with
baseline samples within $1\,\mathrm{mV}$ of an acquisition rail are rejected.
These criteria are frozen and applied to the development and blind acquisitions.

For each mode, the LED threshold is selected on control alone from offsets of
$5$, $15$, $25$, $35$, $45$ and $55\,\mathrm{mV}$ above the event baseline.
An eligible threshold requires at least $95\%$ valid coincident pairs within
$\pm2\,\mathrm{ns}$ of the configured true TOF. Among eligible candidates, the threshold
with the lowest control CTR is retained, with the smaller threshold breaking a tie.
This threshold is then fixed for model selection and blind evaluation.

\subsection{Waveform representation and correction target}
Waveforms are oriented according to detector polarity and limited to the configured
physical amplitude range. The baseline sets the LED crossing level, but the waveform
inputs are not subjected to event-wise baseline subtraction or denoising.
Each detector waveform is cropped relative to its own LED crossing on the native sample
grid and scaled to the interval from zero to one using fixed detector-specific amplitude
limits. Events lacking either finite crossing, the required coincidence, or the complete
input window are excluded.

The UC benchmark uses a short interval from $-1$ to $+2\,\mathrm{ns}$; the FBK
benchmark uses $-1.5$ to $+2\,\mathrm{ns}$. Both use an extended interval from
$-2$ to $+30\,\mathrm{ns}$ relative to the crossing. The UC benchmark evaluates
timing waveforms, whereas the FBK benchmark includes energy and timing waveforms.
The interval is a study setting, rather than a model-dependent crop.
The difference between the LED arrival times, after subtracting the configured true TOF,
is the supervised correction target. If $s_1$ and $s_2$ denote the normalized waveforms and
$h_{\theta}$ the predicted correction, the final residual is
\begin{equation}
 r=(t_{1,\mathrm{LED}}-t_{2,\mathrm{LED}}-\tau)-h_{\theta}(s_1,s_2),
 \label{eq:residual}
\end{equation}
where $\tau$ is the configured true TOF. A correction is therefore learnt for the conventional
timing residual, rather than for an absolute source-position label.

\subsection{Shared and direct estimator architectures}
In a shared formulation, one detector-level scorer is applied independently to both
waveforms:
\begin{equation}
 h_{\theta}(s_1,s_2)=g_{\theta}(s_1)-g_{\theta}(s_2).
 \label{eq:shared}
\end{equation}
Exchanging the detector inputs reverses the correction sign, and identical normalized
pulses yield zero correction. Sharing the scorer removes
arbitrary detector-specific parameter differences while retaining waveform-dependent
corrections. This is an architectural constraint; the adequacy of the constraint for each
board is assessed within that board's data. Alignment to the individual LED crossings
also prevents absolute acquisition times from entering the waveform inputs.

A multilayer perceptron (MLP) implements a fully connected shared scorer.
The configured neural hidden layers use the rectified linear unit (ReLU) activation. A locally
connected MLP uses two temporal layers with position-specific weights, followed by a dense
scorer. Local connectivity retains neighbourhood information without the translation-based
weight sharing of convolution, an established architectural alternative~\cite{elsayed2020}.
Its score difference is passed through an odd, smooth hyperbolic-tangent transformation
with a $500\,\mathrm{ps}$ bound, preserving antisymmetry. A shared one-dimensional CNN uses
one temporal backbone and one dense head for both detectors. Shared linear ridge regression
uses the difference of normalized waveforms without an intercept.

Direct alternatives include an MLP operating on the concatenated pair and linear ridge
regression on the concatenated samples with an intercept. The independent CNN uses two
separately parameterized detector branches and subtracts their scalar outputs; independent
weights do not enforce detector-exchange antisymmetry. The Onishi-inspired paired CNN
fuses the detector dimension in its first convolution, applies further temporal
convolution and uses a dense scalar head~\cite{onishi2022}. Its configured architecture is
a tunable implementation of this correction approach, rather than an assertion of identical
training conditions to the published experiment.

MiniRocket supplies a largely deterministic convolutional feature representation~\cite{minirocket2021}.
The shared variant applies one control-fitted transform and scaling to both detectors,
subtracts their features and fits ridge regression without an intercept. The direct variant
uses a multivariate transform of the paired waveforms and ridge regression with an intercept.
In both cases the transform and scaling are fitted once on control and then frozen; supervised
regression uses development targets only.
For dense MLPs, constant samples are removed using the current fitting population and a
$99\%$ identical-value criterion. Convolutional and locally connected neural architectures
retain the temporal grid; the frozen linear and MiniRocket inputs also retain it.
The estimator inventory is given in Table~\ref{tab:model-inventory}.
\begin{table*}[t]
\centering
\small
\begin{tabularx}{\textwidth}{@{}p{0.14\textwidth}p{0.28\textwidth}X@{}}
\toprule
Formulation & Estimator & Detector processing \\
\midrule
\multirow[t]{5}{*}{Shared} & Dense MLP & One dense scorer for both waveforms; score difference. \\
 & Locally connected MLP & Two position-specific local layers and a dense scorer; smoothly bounded score difference. \\
 & Temporal CNN & One temporal backbone and dense head; score difference. \\
 & Linear ridge & Sample-wise detector difference; no intercept. \\
 & MiniRocket with ridge & One control-fitted transform and scaling; feature difference; no regression intercept. \\
\addlinespace[3pt]
\hdashline
\addlinespace[3pt]
\multirow[t]{5}{*}{Direct} & Dense MLP & Joint dense regression on concatenated waveforms. \\
 & Independent temporal CNN & Two independent backbones and heads; detector-specific score difference. \\
 & Onishi-inspired paired CNN & Detector-fusing first convolution, temporal convolution and dense regression. \\
 & Linear ridge & Concatenated samples with an intercept. \\
 & MiniRocket with ridge & Paired multivariate transform and ridge regression with an intercept. \\
\bottomrule
\end{tabularx}
\caption{Estimator inventory. Shared denotes detector-exchange antisymmetry imposed by a common detector scorer or transform. Direct denotes absence of that imposed constraint, including independently parameterized detector branches.}
\label{tab:model-inventory}
\end{table*}


\subsection{Model development and validation}
Model selection uses only the development acquisition. For each board, mode and window,
the supplied benchmark protocol assigns three deterministic, shuffled folds shared by the
models using outer cross-validation (CV). Every candidate is fitted on the other folds,
and validation CTR, RMSE and the LED reference are computed on the held-out fold.
The candidate spaces in Appendix~\ref{app:spaces} are finite grids in the supplied benchmarks.
Selection minimizes the mean validation CTR. Pruning compares candidates with the LED
reference and completed incumbents on the same evaluated folds using prespecified tolerances;
partially evaluated candidates cannot replace a fully evaluated selection.

After selection, the search fit is discarded and the final estimator is refitted from
scratch using the development population. MLP-based and one-dimensional CNN training use
an RMSE loss, reserving $20\%$ of each fitting population for epoch-level early stopping.
The antisymmetric dense MLP uses stochastic gradient descent (SGD) with Nesterov momentum;
the other such neural models use Adam. Gradient clipping and the configured loss and
gradient stopping criteria are retained. The Onishi-inspired CNN instead uses a mean
squared error (MSE) loss, Adam and a fixed epoch schedule on its full fitting population,
without the internal early-stopping holdout.

The two linear ridge estimators bypass outer CV and candidate pruning. Their regularization
strength is selected by internal leave-one-out CV using MSE, and the final regression is
fitted on the full development population. Internal ridge MSE is not an outer-validation
CTR and will not be entered as one in the result tables. MiniRocket ridge regression remains
within the outer-CV protocol. One configured random seed determines splitting, model
training and subsequent resampling; multiple-seed replicas are not part of the benchmark.

\subsection{Blind evaluation and temporal sensitivity}
The final model is applied to the independent blind acquisition with all preprocessing,
threshold and model choices frozen. Predicted corrections are limited to an absolute
magnitude of $2000\,\mathrm{ps}$ before residual evaluation; the locally connected model
also retains its narrower smooth bound. LED and corrected residuals use the same events.
CTR and RMSE uncertainties are estimated by paired blind-event bootstrap resampling,
without retraining. The supplied settings use $1000$ draws for UC and $100$ for FBK. These uncertainties describe
event-level variation conditional on the fitted model, whereas validation dispersion is
fold-to-fold variability. Within-board differences between models are also evaluated by
matching blind-event identities and resampling the matched events jointly.

Explainable artificial intelligence (XAI) is based on temporal occlusion. Groups of eight
retained samples are replaced by interpolation between adjacent samples, and the mean
absolute change in the predicted correction is measured over a reproducible subset of at
most $4096$ blind events. Shared estimators perturb both detector inputs together, producing
one temporal importance profile. Direct estimators perturb one detector at a time, producing
two channel-specific profiles. For display, importance is aggregated into $1\,\mathrm{ns}$
bins and normalized; the direct channels use a common maximum. These maps quantify model
sensitivity to the defined perturbation, rather than a causal attribution. XAI is not used
for hyperparameter selection.
